\documentclass[11pt]{article}
\usepackage[margin=1in]{geometry} 
\usepackage[parfill]{parskip}% Begin paragraphs with an empty line rather than an indent
\usepackage[largesc]{newtxtext}
\usepackage[scaled=1.1]{nimbusmononarrow}
\usepackage[varbb]{newtxmath}
\title{{\tt\Large pxtxalfa}\\Math Alphabets Derived From {\tt\Large newpx} and {\tt\Large newtx}}
\author{Michael Sharpe}
%\email{msharpe at ucsd dot edu}

\begin{document}
\maketitle
\section{Overview}
The {\tt txfonts} and {\tt pxfonts} packages, both created by Young Ryu but no longer under  active development, provide fairly complete typesetting environments based on the Times and Palatino text font families respectively. Other packages (eg, {\tt txgreeks}, providing the option of upright or slanted Greek letters) extend the range of coverage of its macros. 

These packages contain some interesting math alphabets. The script alphabet glyphs (upper case only) seem to be identical to those in {\tt Mathematica5}, but the Fraktur font common to both packages is. as far as I can tell,  distinct from the Fraktur of other major math font packages, and worthy of note. Blackboard bold comes in two different versions in {\tt txfonts} (openface and double-struck) and in yet another double-struck version in {\tt pxfonts}. The double-struck alphabets are similar in overall style to those in {\tt mathpazo} and {\tt Mathematica7}, with stems a mix of double-struck, regular weight  and solid bold.


The original plan was to provide virtual fonts for all these alphabets, plus packages that allow them to be used in stand-alone fashion and as part of the \textsf{mathalfa} package.

In the decade since this package was first released, some changes have been made to the original alphabets, most notably:
\begin{itemize}
\item
the Fraktur fonts have been modified so that they now match the weight of Times-like fonts rather than lighter fonts such as Computer Modern;
\item The secondary Blackboard Bold font provided with {\tt newtx}, with uppercase letters of the form $\mathbb{ABCDEXYZ}$ has been extended with full lowercase plus {\tt dotlessi} and {\tt dotlessj}, like $\mathbb{abcdexyz}$.
\end{itemize}


The package contains the following files: 

Virtual fonts ({\tt.tfm} and {\tt.vf}):

\begin{tabular}{lll}
{\tt txr-cal}&Regular weight calligraphic from {\tt txfonts} and {\tt pxfonts}.\\
{\tt txb-cal}&Bold weight calligraphic from {\tt txfonts} and {\tt pxfonts}.\\
{\tt txr-frak}&Regular weight fraktur from {\tt txfonts} and {\tt pxfonts}.\\
{\tt txb-frak}&Bold weight fraktur from {\tt txfonts} and {\tt pxfonts}.\\
{\tt txr-of}&Regular weight openface from {\tt txfonts}.\\
{\tt txb-of}&Bold weight openface from {\tt txfonts}.\\
{\tt txr-ds}&Regular weight double-struck from {\tt txfonts}.\\
{\tt pxr-ds}&Regular weight double-struck from {\tt pxfonts}.\\
{\tt pxb-ds}&Bold weight double-struck from {\tt pxfonts}.
\end{tabular}

Font definition ({\tt.fd}) files:

\begin{tabular}{lll}
{\tt utx-cal.fd}&Regular and bold weights, calligraphic.\\
{\tt ot1tx-frak.fd}&Regular and bold weights, fraktur.\\
{\tt utx-of.fd}&Regular and bold weights, openface.\\
{\tt ot1tx-ds.fd}&Regular weight double-struck from {\tt txfonts}.\\
{\tt upx-ds.fd}&Regular and bold weights, double-struck from {\tt pxfonts}.
\end{tabular}

Style files:

\begin{tabular}{lll}
{\tt pxtx-cal.sty}&Load regular and bold weights, calligraphic.\\
{\tt pxtx-frak.sty}&Load regular and bold weights, fraktur.\\
{\tt tx-of.sty}&Load regular and bold weights, openface.\\
{\tt tx-ds.sty}&Load regular weight double-struck from {\tt txfonts}.\\
{\tt px-ds.sty}&Load regular and bold weights, double-struck from {\tt pxfonts}.\\
\end{tabular}
(Only {\tt pxtx-frak.sty} and {\tt tx-ds.sty} have been modified since the the original versions.)
When loaded following loading other math sty files, these will replace one of the math alphabets. For example, {\tt tx-ds.sty} will redefine \verb|\mathbb| and its associated special characters, like \verb|\bbdotlessi|, to use the double-struck glyphs from {\tt newtx}.
 
\section{The interesting font files}
The files ({\tt.afm} and {\tt.pfb}) with glyphs of interest are:
\begin{verbatim}
txmiaX, txbmiaX---Fraktur (UC, lc) and Double-Struck (regular weight only)
txsy, txbsy---Calligraphic (UC)
txsyb, txbsyb---Openface (UC)
pxsyb, pxbsyb---Double-Struck (UC)
\end{verbatim}

 This package depends on {\tt txfonts} and {\tt pxfonts}. It will not function unless the map files {\tt txfonts.map} and {\tt pxfonts.map} are enabled. This is the default in \TeX\ Live installations.


On the other hand, the metrics for the math alphabets in this collection have been adjusted and do not have the problems of the originals. This is a matter of personal taste, and may not suit yours. Sorry---there is no way to allow simple user-configured settings for these parameters.

The easiest way to use the fonts in this package is {\tt mathalpha}, \textsc{aka} {\tt mathalfa}, the latest version of which builds in support for these alphabets. For font samples, see the documentation for that package.

\end{document}  

